\documentclass[10pt,a4paper]{amsart}
\usepackage[margin=19mm]{geometry}
\usepackage{amsmath,amssymb,mathtools}
\usepackage[scr=rsfs,cal=euler]{mathalfa}
\usepackage{xfrac}
\usepackage[unicode,hidelinks,bookmarks=false]{hyperref}
\newcommand{\ie}{i.e.\ }
\newcommand{\cf}{cf.\ }
\newcommand{\resp}{respectively\ }
\newcommand{\Subsection}{\subsection}
\providecommand{\nobreakdash}{\nobreak\hbox{-}}
\setlength{\emergencystretch}{2em}
\multlinegap0pt
\usepackage{amssymb}
\usepackage[shortlabels]{enumitem}
\usepackage{tikz}
\newtheorem{lemma}{Lemma}
\newtheorem{claim}{Claim}
\usepackage{cleveref}
\crefname{lemma}{Lemma}{Lemmas}
\crefname{claim}{Claim}{Claims}
\newcommand{\cF}{\mathcal{F}}
\title{Cross-free families have linear size}
\author{Istv\'an Tomon}
\date{Source: arXiv:2603.05443v3 (CC BY 4.0)}
\begin{document}
\maketitle

\noindent\emph{Source and licence.} Cross-free families have linear size. Version arXiv:2603.05443v3 (CC BY 4.0), distributed by arXiv under the Creative Commons Attribution 4.0 International licence, http://creativecommons.org/licenses/by/4.0/, as recorded in the arXiv licence metadata for this exact version and shown on the abstract page https://arxiv.org/abs/2603.05443v3. Full licence notice: \texttt{licenses/arxiv-2603.05443-CC-BY-4.0.txt}.\par\smallskip

\noindent\textit{Editorial scope.} Counted unit: the article's lemma lemma (source0.tex lines 364--371). The article's complete original proof of it is delivered with it (source0.tex lines 372--409, one \texttt{proof} environment of 38 lines). No mathematics is written, repaired or re-derived: the statement and the proof are the article's own lines, in the article's own notation, and no retained line is substituted or re-flowed.  Retained uncounted as the source-local context the proof needs: lines 145--147; lines 333--335.  Nothing was omitted from any retained span, so the package carries no omission record.  No retained line contains a citation, so the wrapper carries no bibliography and no bibliography entry.  No retained line contains a figure environment, an image-inclusion command or drawing code, so no image is delivered and the package declares no asset dependency.  Editorial formatting only, in addition: an AMS \texttt{amsart} preamble that restates the article's own declarations and macros used by the retained lines, namely the environments \texttt{lemma}, \texttt{claim} on separate counters, together with the standard packages those lines need.  The retained environments print in this document as Lemma 1, Claim 1 in order of appearance, whereas the article prints the same items with the numbering of its own full document.  The complete native source of the article as distributed in the arXiv e-print for this exact version is bundled unchanged as \texttt{sources/195803/source0.tex}.
\begin{lemma}\label{lemma:dilworth}
Let $\cF\subset 2^X$ be a $k$-wise laminar and $k$-wise intersecting family. Then $\mathcal{F}$ can be partitioned into $k-1$ chains.
\end{lemma}

\begin{lemma}\label{lemma:T}
$T$ satisfies \textbf{T6}-\textbf{T9}. Moreover, \textbf{T6}-\textbf{T8} do not use \textbf{T5}.
\end{lemma}

 \begin{lemma}
 For $\ell=0,1,\dots,k$, there exists $I_{\ell}\subset I$ of size $|I_{\ell}|\geq |I|-\ell h n$, and for every $i\in I_{\ell}$, there exists a cross-support tree $T_{\ell,i}$ such that 
\begin{itemize}
    \item[(1)] the height of $T_{\ell,i}$ is $\ell$,
    \item[(2)] the root of $T_{\ell,i}$ is $i$,
    \item[(3)] every non-leaf vertex of $T_{\ell,i}$ has at least $h/(2k)^{k\ell}$ children.
\end{itemize}
\end{lemma}

\begin{proof}
We proceed by induction on $\ell$. The case $\ell=0$ follows by taking $I_0=I$, and defining $T_{0,i}$ to be the single vertex tree, whose root is $i$. Now assume that $\ell>0$ and that the statement holds for $\ell-1$ instead of $\ell$.

For every $i\in I_{\ell-1}$, let $Y_i=X_i$ if $\ell=1$, and let $Y_i$ be the set of labels on the edges connecting the root $i$ of $T_{\ell-1,i}$ to its children if $\ell\geq 2$. Then $Y_i\subseteq X_i$ by \textbf{T2} and $|Y_i|\geq h/(2k)^{k(\ell-1)}$. Let $Z_i$ be the second half of $Y_i$ with respect to the ordering $\prec$ of $X$, so $Z_i\subset Y_i$ and $|Z_i|\geq |Y_i|/2$. For each $x\in X$, let 
$$J_x=\{i\in I_{\ell-1}:x\in Z_i\}.$$
By \textbf{C1}, the family $\{C_i(x):i\in J_x\}$ is a chain. Let $J^{\text{top}}_x\subset J_x$ be the set of indices corresponding to the $h$ largest members of this chain, or $J^{\text{top}}_x=J_x$ if $|J_x|<h$. Define 
$$I_{\ell}:=I_{\ell-1}\setminus \bigcup_{x\in X}J_x^{\text{top}},$$
then $|I_{\ell}|\geq |I_{\ell-1}|-hn\geq |I|-h\ell n$.

For every $i\in I_{\ell}$, we define the desired tree $T_{\ell,i}$ as follows. For every $x\in Z_i$, there are at least $h$ indices $j\in I_{\ell-1}$ such that $x\in Z_j$ and $C_i(x)\subset C_j(x)$. This is true because $i\in J_x\cap I_{\ell}$, which means that $|J_x^{\text{top}}|=h$ and every $j\in J_x^{\text{top}}$ works. Therefore, for each $x\in Z_i$, we can pick an index $j_x\in J_x^{\text{top}}$ such that $j_x\neq j_y$ for $x\neq y\in Z_i$. If $\ell=1$, set $T'_x=T_{0,j_x}$. If $\ell\geq 2$, then for each tree $T_{\ell-1,j_x}$, remove all children of the root $j_x$ (and their descendants), whose joining edge is labeled by an element of $X$ that is larger than $x$ in the ordering $\prec$, and let $T'_{x}$ be the resulting tree. By \textbf{T9}, $T'_x$ is also a cross-support tree. Recalling that $x\in Z_{j_x}$, $x$ is in the second half of $Y_{j_x}$ with respect to $\prec$, so the root of $T_{x}'$ has degree at least $h/(2(2k)^{k(\ell-1)})>h/(2k)^{k\ell}$. To summarize, we have the following properties of $T_x'$:
\begin{itemize}
    \item $T_{x}'$ is a cross-support tree with root $j_x$,
    \item if $\ell\geq 2$, the leftmost edge from the root is labeled with  $x$,
    \item $C_i(x)\subset C_{j_x}(x)$,
    \item every non-leaf vertex has at least $h/(2k)^{k\ell}$ children. 
\end{itemize}

Now define the tree $T'$ by taking the disjoint union of the trees $T_{x}'$ for $x\in Z_i$, and joining the root $j_x$ of $T_x'$ to the new root $i$. The trees $T_{x}'$, $x\in Z_i$, are drawn from left to right with respect to the reverse of the order $\prec$ on $X$.  We label the edge $ij_x$ by $x$. Clearly, the highlighted properties of the trees $T_x'$ ensure that $T'$ satisfies \textbf{T1}-\textbf{T4}, which then imply \textbf{T6}-\textbf{T8} by \Cref{lemma:T}. However, \textbf{T5} does not necessarily hold, so we remove some further children of the root of $T'$.

 For $d=0,\dots,\ell-1$, let $i_{d,x}$ be the leftmost vertex of $T'_x$ at depth $d$. The family 
$$\{S_{i_{d,x}}:x\in Z_i\}$$ is $k$-wise intersecting by \textbf{T7}. Therefore, applying \Cref{lemma:dilworth} in a total of $\ell$ times, there exists $Q\subset Z_i$ of size at least 
$$|Q|\geq \frac{|Z_i|}{(k-1)^{\ell}}\geq \frac{|Y_i|}{2(k-1)^{\ell}}\geq \frac{h/(2k)^{(\ell-1) k}}{2(k-1)^{\ell}}>\frac{h}{(2k)^{k\ell}}$$
such that $$D_d=\{S_{i_{d,x}}:x\in Q\}$$ is a chain for every $d=0,\dots,\ell-1$. Let $T=T_{\ell,i}$ be the tree we get by keeping only those subtrees $T'_x$ of $T'$ for which $x\in Q$. 

\begin{claim}
$T$ satisfies \textbf{T5}.
\end{claim}

\begin{proof}
First, we show that if $x,y\in Q$ and $y\prec x$, then $S_{i_{d,y}}\subset S_{i_{d,x}}$. As $D_d$ is a chain, we have that $S_{i_{d,x}}$ and $S_{i_{d,y}}$ are comparable. As $i_{d,x}$ is a leftmost vertex in $T_x'=T[j_x]$, we have $C_{i}(x)\subset S_{j_x}\subset S_{i_{d,x}}$ by \textbf{T6}. But $y\prec x$, so $C_{i}(y)\cup\{y\}\subseteq C_{i}(x)$ by \textbf{C2}, and in particular $y\in C_i(x)$. As $i_{d,y}$ is a leftmost vertex in $T_{y}'$, we have $\phi(i_{d,y})=y$ by \textbf{T6}, so $S_{i_{d,y}}=C_{i_{d,y}}(y)$. In particular, $y\not\in S_{i_{d,y}}$. Thus, as $S_{i_{d,x}}$ and $S_{i_{d,y}}$ are comparable, we must have $S_{i_{d,y}}\subset S_{i_{d,x}}$.

In order to prove that $T$ satisfies \textbf{T5}, we need to verify the following. Let $u,u'\in V(T)$ be distinct vertices at the same depth such that $u$ is to the left of $u'$. Let $a$ be their common ancestor of maximum depth, and let $v\neq v'$ be the children of $a$, where $v$ is the ancestor of $u$ (or $v=u$), and $v'$ is the ancestor of $u'$ (or $v'=u'$). If $u$ is a leftmost vertex in $T[v]$, then $S_{u'}\subseteq S_u$.

If the common ancestor $a$ is not the root $i$, then this is true using that $u$ and $u'$ lie in some cross-support subtree $T'_x$, $x\in Q$. Hence, assume that $a=i$, then $v=j_x$ and $v'=j_y$ for some $x,y\in Q$, $y\prec x$. Also, $u$ is a leftmost vertex in $T_x'$, so $u=i_{d,x}$ for some $d$. We have $S_{u'}\subset S_{i_{d,y}}$ using that $u',i_{d,y}$ are vertices of the cross-support tree $T_y'$. But then $S_{u'}\subset S_{i_{d,y}}\subset S_{i_{d,x}}=S_u$, verifying the required claim.
\end{proof}

In conclusion, $|I_{\ell}|\geq |I|-\ell h n$, and for every $i\in I_{\ell}$, the tree $T=T_{\ell,i}$ is a cross-support tree of height $\ell$ in which every non-leaf vertex has at least $h/(2k)^{k\ell}$ children. This finishes the proof.
\end{proof}

\end{document}
